\scanpage{71}
\mspage{F.}
% archive foliation: 70
\margin{\uline{Pic birationnel}}
\uline{Théorème} Soient $f:X\to Y$ un morphisme dominant
\inserted{birationnel}
de préschémas intègres normaux \deleted{et} propres sur $k$,
avec $Y$ \deleted{factoriel}\inserted{\deleted{\unsure{loc.}}\uline{régulier}}.
Alors le morphisme
\[
\underline{\mathrm{Pic}}_Y^\tau\longrightarrow
\underline{\mathrm{Pic}}_X^\tau
\]
est un \uline{isomorphisme}.

\margin{$X$ normal ``sur $k$'' ?}
\margin{peut-on \deleted{\unsure{supposer $Y$ régulier??}}
\deleted{\unsure{connexe}} \unsure{sur $k$} ?}
\margin{N.B. factoriel ne suffit pas $=$ \unsure{est des pts de ces}
\unsure{non-réduits, cf côtés} \ldots}

C'est un \uline{monomorphisme} grâce au fait que
$f_*(\underline O_X)=\underline O_Y$ (universellement sur $k$ !),
donc une immersion fermée puisque un hom de groupes
\unsure{-- en principe} propre.
Montrons que c'est \uline{surjectif}.
Soit $U$ l'ouv de $Y$ où $f^{-1}$ est défini,
comme $Y$ est factoriel, on a
\[
\mathrm{Pic}(X)\xrightarrow{\simeq}\mathrm{Pic}(U),
\]
d'où, utilisant
\diagram{\begin{tikzcd}
&X\arrow[d,"f"]\\
U\arrow[ur,"g"]\arrow[r,"i"']&Y
\end{tikzcd}},
le diagramme commutatif
\diagram{\begin{tikzcd}
&\mathrm{Pic}(X)\arrow[dl,leftrightarrow,"g^*"']\\
\mathrm{Pic}(U)&\mathrm{Pic}(Y)\arrow[l,"\simeq","i^*"']\arrow[u,"f^*"']
\end{tikzcd}}
qui montre que $\mathrm{Pic}(Y)\to\mathrm{Pic}(X)$ est un hom
sur un facteur direct, l'image étant formée des $\xi$ telles que
\deleted{$g^*(\xi)=0$.}
$\xi=f^*i^{*-1}g^*(\xi)$.
Or si $\xi$ est \deleted{\unsure{même}} $\tau$-équivalente à $0$,
\inserted{donc $\eta=m\xi$ est alg. équiv. à $0$,}
\deleted{\unsure{alors}}
$g^*(\unsure{\eta|\xi})$ est
\deleted{$\tau$-équivalente}\inserted{alg. équiv.} à $0$,
donc $i^{*-1}g^*(\unsure{\eta|\xi})\deleted{\unsure{même}}$ aussi,
donc $mf^*i^{*-1}g^*(\xi)$
\inserted{$=mf^*(i^{*-1}g^*(\xi))$}
\deleted{\unsure{aussi}}.
D'autre part, on voit que
\inserted{$\xi-f^*i^{*-1}g^*(\xi)$}
\deleted{\unsure{le diviseur exceptionnel}}
est représenté par un

\margin{Théorie des \unsure{classes} des cycles !}
\margin{$Y$ projectif ?? ou \unsure{proj.} sur $k$}
